%% ============================================================
%% 分片 03：原文第 2 节「${\EUN}=4$ Super Yang-Mills Theory And $S$-Duality」（原文 glshort.tex 起于第 642 行）
%% 由 tools/harvmac2tex.py 从 arXiv 源码（hep-th/0604151v3）自动生成。
%% 只译英文散文；公式、编号、\label、\cite 一律不改。
%% ============================================================

\subsection{$S$-Duality}

\label{sec:reviewsd}

A review of $S$-duality is unfortunately beyond our scope in this
paper.  We will just mention a few relevant facts.

Since its imaginary part is positive, $\tau=\theta/2\pi +4\pi
i/e^2$ defines a point in the upper half plane ${\EUH}$.  The
group $SL(2,\Bbb{R})$ acts on ${\EUH}$ in the standard fashion
$\tau\to (a\tau+b)/(c\tau+d)$, with $ad-bc=1$.\footnote{The group that
acts faithfully on ${\EUH}$ is the quotient $PSL(2,\Bbb{R})$, but
in application to four-dimensional gauge theory, one really needs
the double cover.} The transformation $T:\tau\to\tau+1$ is simply
a $2\pi$ shift of the angle $\theta$ and thus a classical symmetry
of the theory on $\Bbb{R}^4$, for any gauge group $G$. The
$S$-duality conjecture asserts that there exists an additional
quantum symmetry that inverts $\tau$, exchanges $G$ with
$^L\negthinspace G$, and exchanges electric and magnetic charges.
Moreover, this symmetry, which we will call $S$, combines with the
classical symmetry $T:\tau\to\tau+1$ to generate an infinite
discrete subgroup $\Gamma$ of $SL(2,\Bbb{R})$.

The most familiar case is the case that $G$ is simply-laced.  Then
% 式 (2.16)
\setcounter{equation}{15}
\begin{equation}\label{eq:bivolo}
S=\left(\begin{matrix} 0 & 1\\ -1 & 0\\\end{matrix}\right)
\end{equation}
 acts as
$\tau\to -1/\tau$, and together with $T$ generates the group
$SL(2,\Bbb{Z})$.

If $G$ is not simply-laced, then the $S$-transformation is not
$\tau\to -1/\tau$.  Rather, it is $\tau\to -1/n_{\frak g}\tau$,
where $n_{\frak g}$ is $2$ for $F_4$ and 3 for $G_2$.  This
transformation can be achieved by the $SL(2,\Bbb{R})$
transformation % [WARN] 公式 \eqn\eqref{eq:bivolo} 无编号或重复
\begin{equation*}
\tag*{(2.17)}
S=\left(\begin{matrix} 0 & 1/\sqrt {n_{\frak
g}}\\ -\sqrt {n_{\frak g}} & 0\\\end{matrix}\right).
\end{equation*}
 This presence of a
factor of 2 or 3 in the action of $S$ on $\tau$ can be seen
\cite{ref46,ref47} by examining the BPS mass formulas for
electric and magnetic charges
 and reflects the relation between roots and
coroots for these groups. It also can be extracted from a
string-theoretic explanation of $S$-duality for non-simply-laced
groups \cite{ref48}. The factor of 2 or 3 means that the duality
groups for $G_2$ or $F_4$ are not simply $SL(2,\Bbb{Z})$, but
certain infinite discrete subgroups of $SL(2,\Bbb{R})$ that are
known as Hecke groups.

The remaining simple Lie groups $Sp(k)$ and $Spin(2k+1)$ (and
their respective quotients by $\Bbb{Z}_2$) require some special
comment, because these are the only simple Lie groups such that
$G$ and $^L\negthinspace G$ have non-isomorphic Lie algebras.  The
Yang-Mills Lagrangian and therefore the definition of the $\tau$
parameter depend only on the Lie algebra $\frak g$ of the gauge
group. Hence, whenever $G$ and $^L\negthinspace G$ have the same
Lie algebra, we can discuss how the $S$-transformation acts on
$\tau$ without distinguishing whether we have in mind $\tau$ of a
theory with gauge group $G$ or $\tau$ of a theory with gauge group
$^L\negthinspace G$. This indeed is what we have implicitly done
so far.

For the pair $Sp(k)$ and $Spin(2k+1)$, however, there is no
equally direct way to compare the two $\tau$ parameters.  Hence,
one may introduce separate gauge coupling parameters, say $\tau$
for $Sp(k)$ or $Sp(k)/\Bbb{Z}_2$ gauge theory, and $\tau'$ for
$Spin(2k+1)$ or $Spin(2k+1)/\Bbb{Z}_2=SO(2k+1)$ gauge theory. If
one normalizes the respective $\tau$ parameters so that the
respective $T$-transformations act by $\tau\to \tau+1$ and
$\tau'\to\tau'+1$, then $S$ acts by $\tau=-1/2\tau'$, $\tau'=
-1/2\tau$, just as for $F_4$.  (Thus we set $n_{\frak g}=2$ for
$Spin(2k+1)$ and $Sp(k)$.)

This normalization is natural, since it leads to the most uniform
gauge theory formulas for arbitrary gauge groups. A slight
complication is that there is for the same groups a second
normalization that one might also consider natural.  A useful
string theory realization of the $Sp(k)/Spin(2k+1)$ duality,
involving orientifold threeplanes \cite{ref49}, actually
motivates a different normalization.  In this normalization, the
$Sp(k)$ theory is described by a coupling parameter
$\tilde\tau=2\tau$, so that $S$ acts simply by $\tilde \tau =
-1/\tau'$, but instead $T$ acts by $\tilde\tau\to \tilde \tau+2$.
(In figure 3 of \cite{ref49}, the $Sp(k)$ theory appears twice,
precisely because  the $Sp(k)$ coupling parameter $\tilde\tau$ is
normalized in that paper so that the $T$-symmetry acts by
$\tilde\tau\to\tilde\tau+2$.)

A final comment on this is that the assertion that Montonen-Olive
duality exchanges $G$ and $^L\negthinspace G$ is not quite the
whole story.  On $\Bbb{R}^4$, this is an adequate description, but
on a general four-manifold, the full story is somewhat more
elaborate.  A $G$-bundle on a four-manifold has a characteristic
class $\xi\in H^2(M,\pi_1(G))$, studied in this paper in section
第~\ref{sec:centertop}~节. A path integral can be defined for each value of
$\xi$, and the resulting partition functions $Z_\xi$ transform as
a unitary representation of the duality group
\cite{ref45,ref26}.

How much of this is important for the geometric Langlands program?
The basic geometric Langlands duality is the transformation $S$
that acts as $\tau\to -1/n_{\frak g}\tau$.  It acts on the
canonical parameter $\Psi$ (introduced in section 第~\ref{sec:canonpar}~节) by
$\Psi\to -1/n_{\frak g}\Psi$.  The basic geometric Langlands
duality involves comparing $\Psi=\infty$ to $\Psi=0$.  These are
exchanged by $S$ regardless of the value of $n_{\frak g}$.  So the
value of $n_{\frak g}$ is not very important for the basic
geometric Langlands duality.  In section 第~\ref{sec:genotics}~节, we come to a
generalization of the geometric Langlands duality to arbitrary
$\Psi$.  Here the precise duality group is important and,
therefore, the value of $n_{\frak g}$ does play a role.

\par\medskip{\it Transformation of Supersymmetries}

There is one question about Montonen-Olive duality that actually
will play a bigger role in this paper: How does it act on the
supersymmetries?

 We cannot the answer this question by inspection
of the classical Lagrangian, because, apart from the subgroup
generated by $\tau\to \tau+1$, $\Gamma$ does not consist of
symmetries of the classical theory. This after all is what makes
$S$-duality interesting.  So obtaining the answer will require a
more subtle reasoning.

Before determining the answer, let us ask to what extent the
answer is unique.  Consider an element $\gamma=\left(\begin{matrix}a & b
\\ c & d \\\end{matrix}\right)$ of the duality group $\Gamma$ generated by
$S$ and $T$. Such an element acts on $\tau$ by $\tau\to
(a\tau+b)/(c\tau+d)$ and on the supersymmetry algebra by an
automorphism. This automorphism is not uniquely defined, since it
could be combined with a symmetry of the ${\EUN}=4$ super
Yang-Mills theory (at a fixed value of $\tau$). An important
simplification is that, according to the Montonen-Olive
conjecture, $\gamma$ commutes with the Poincar\'e group. Moreover,
one can define it to commute with  the global $R$-symmetry group
$SU(4)_{{\cal R}}$.\footnote{To show this, we first observe that
conjugation by $\gamma$ generates an automorphism of $SU(4)_{{\cal
R}}$.  This automorphism is necessarily inner, as the classical
theory has no symmetry that acts trivially on spacetime and by an
outer automorphism of $SU(4)_{{\cal R}}$. Finally, given that
$\gamma$ generates an inner automorphism of $SU(4)_{{\cal R}}$, we
can combine it with an $SU(4)_{{\cal R}}$ element that generates
the inverse automorphism to get a duality symmetry that commutes
with $SU(4)_{{\cal R}}$. } Combining these facts, it follows that
$\gamma$ acts as a scalar multiplication $\exp(i\hat\phi)$ on the
supersymmetries that transform as the $\bar{\bf 4}$ (that is, the
$\bar Q_{AX}$) and as $\exp(-i\hat\phi)$ on those that transform
as ${\bf 4}$ (the $ Q_{\dot A}^Y$). Moreover, $\hat\phi$ must be
real to preserve the real structure of the algebra (with respect
to which, in Lorentz signature, $P$ is hermitian and $\bar Q$ is
the hermitian adjoint of $Q$).  We will call these symmetries
$U(1)$ chiral rotations. The action of $\gamma$ is defined up to
an element of the center of $SU(4)_{{\cal R}}$. The center is
generated by an element ${\EUI}$ that acts as $i$ on the ${\bf 4}$
of $SU(4)_{{\cal R}}$ and $-i$ on the $\bf{ \bar 4}$. Thus
$\hat\phi(\gamma)$ is defined up to
$\hat\phi(\gamma)\to\hat\phi(\gamma)+\pi/2$.

To determine $\hat\phi(\gamma)$, we can compute the action of
$\gamma$ in any convenient state. We choose to perform the
computation on the Coulomb branch of the theory, where the gauge
group is broken to $U(1)^r$ and the supersymmetry algebra is
centrally extended by electric and magnetic charges.  For an
abelian gauge group such as $U(1)$ or $U(1)^r$, one can calculate
everything explicitly and thereby determine how $\gamma$ acts on
the supersymmetries.  Or one can use the realization of ${\cal
N}=4$ super Yang-Mills theory as the gauge theory on a $D3$-brane
of Type IIB superstring theory; this gives a geometrical way to
determine the action of $\gamma$ on supercharges for $G=U(1)$ (and
this approach also extends directly to classical groups such as
$U(N)$).

We will follow the alternative approach of determining the result
by examining the mass formula for BPS states. To simplify
notation, we focus on an ${\EUN}=2$ subalgebra of the
supersymmetry algebra, which has a pair of right-handed
supercharges $Q_{\dot A}^i$, $i=1,2$ and a single central charge
$Z$. They satisfy % 式 (2.18)
\setcounter{equation}{17}
\begin{equation}\label{eq:hilj}
\{Q_{\dot A}^i,Q_{\dot
B}^j\}=\eps_{\dot A\dot B}\eps^{ij} Z.
\end{equation}
 The complex scalar which
is the ${\EUN}=2$ superpartner of the massless gauge fields takes
values in the Cartan subalgebra ${\frak t}$ of the Lie algebra
${\frak g}$. We normalize this scalar so that its kinetic term is
$\tau$-independent and denote its expectation value by
$\vec{\phi}$.  The electric and magnetic charges, denoted by
$\vec{\eurm n}$, $\vec{\eurm m},$ take values in the weight and
coweight lattices of $\frak g$. The vectors $\vec{\phi},\vec{\eurm
n},\vec{\eurm m}$ are defined up to an action of the Weyl group.

To make our conventions clear, we will give the explicit
definition of $\eurm n$ and $\eurm m$ for $G=U(1)$.  Physically,
one usually describes $U(1)$ gauge theory with a real connection
$A'$ of curvature $F'=dA'$.  We write Minkowski spacetime as
$\Bbb{R}^{1,3}=\Bbb{R}\times \Bbb{R}^3$ where the first factor
parametrizes time and the second parametrizes space.  To agree
with Maxwell's equations (see, for example, p. 42 of
\cite{ref51}), one defines the electric and magnetic fields by
$E_i=F'{}^{0i}$ and $B_i=\half\epsilon_{ijk}F'_{jk}$ (here
$\epsilon_{ijk}$ is the antisymmetric tensor on $\Bbb{R}^3$ with
$\epsilon_{123}=1$). Then $\eurm n$ and $\eurm m$ are defined by
% 式 (2.19)
\setcounter{equation}{18}
\begin{equation}\label{eq:eltro}
\begin{aligned}{\eurm n} & = {1\over 4\pi}\int_Y d\Omega\,
\vec n\cdot \vec E\\ {\eurm m}&={1\over 4\pi}\int_Yd\Omega\,\vec
n\cdot \vec B\\\end{aligned}
\end{equation}
 where $Y$ is a large sphere at infinity in
$\Bbb{R}^3$ with volume-form $d\Omega$ and normal vector $\vec n$.
In terms of the antihermitian gauge field $A$ related to $A'$ by
$A=-iA'$ and $F=-iF'$ (recall eqns. \eqref{eq:boloko}\ and \eqref{eq:obolko}), we have
% 式 (2.20)
\setcounter{equation}{19}
\begin{equation}\label{eq:feltro}
\begin{aligned}{\eurm n} & = {i\over 4\pi}\int_Y d\Omega\,
n^i F^{0i}\\ {\eurm m}&={i\over 4\pi}\int_Yd\Omega\,\half
\epsilon_{ijk}n^i F^{jk}.\\\end{aligned}
\end{equation}
 In terms of differential forms, the
second formula is % 式 (2.21)
\setcounter{equation}{20}
\begin{equation}\label{eq:ofeltro}
{\eurm m}={i\over 2\pi}\int_YF
=\int_Y c_1({\cal L}),
\end{equation}
 where ${\cal L}$ is a line bundle on which
$A$ is a connection.

The central charge $Z$ is given by % 式 (2.22)
\setcounter{equation}{21}
\begin{equation}\label{eq:pikf}
Z=\sqrt{2\over {\rm
Im}\,\tau}\,\left(\vec{\eurm m}~\vec{\eurm n}\right)\cdot
\left(\begin{matrix}\tau\vec \phi \\\vec \phi \\\end{matrix}\right)=\sqrt {2\over
{\rm Im}\,\tau}\,\vec \phi\cdot (\vec{\eurm n}+\tau \vec{\eurm m})
\end{equation}
  For simply-laced $G$, the low energy abelian gauge theory on
 the Coulomb branch admits an $SL(2,\Bbb{Z})$ group of duality
 symmetries such that a general element
$\eurm M=\left(\begin{matrix}a&b\\ c&d\\\end{matrix}\right)$, which  acts on
$\tau$ by $\tau\to (a\tau+b)/(c\tau+d)$, acts on $\vec{\phi},\,
\vec{{\eurm m}},\, \vec{{\eurm n}} $ by
% 式 (2.23)
\setcounter{equation}{22}
\begin{equation}\label{eq:xilj}
\begin{aligned}\vec{\phi} & \to
\vec{\phi}\\\left(\begin{matrix}\vec{{\eurm m}} & \vec{{\eurm
n}}\end{matrix}\right)&\to \left(\begin{matrix}\vec{{\eurm m}} & \vec{{\eurm
n}}\end{matrix}\right){\eurm M}^{-1}=\left(\begin{matrix} \vec{{\eurm m}} &
\vec{{\eurm n}}\end{matrix}\right)\left(\begin{matrix}d & -b  \\ -c & a
\\\end{matrix}\right).\end{aligned}
\end{equation}
 The $S$-duality conjecture asserts that this
symmetry of the low energy theory actually extends to a symmetry
of the full theory, necessarily acting in the same way. That the
generator $T:\tau\to\tau+1$ of the duality group $\Gamma$ extends
to the full theory is clear (since it is just a $2\pi$ shift in
the $\theta$ angle), so the conjecture really is that $S:\tau\to
-1/\tau$ similarly extends to the full theory. The action of $T$
claimed in \eqref{eq:xilj}\ can be seen in the full theory by a direct
computation \cite{ref50}.  This direct computation is valid for
all $G$, simply-laced or not; a variant of it will be presented in
section 第~\ref{sec:tophooft}~节. From \eqref{eq:pikf}\ and \eqref{eq:xilj}, one deduces the
transformation of the central charge: % 式 (2.24)
\setcounter{equation}{23}
\begin{equation}\label{eq:mujk}
Z\to
{|c\tau+d|\over {c\tau+d}}\, Z.
\end{equation}
 Then the centrally-extended
supersymmetry algebra implies that the right-handed
supersymmetries get multiplied by
% 式 (2.25)
\setcounter{equation}{24}
\begin{equation}\label{eq:tokio}
\exp(-i\hat\phi(\gamma))=\left({|c\tau+d|\over
{c\tau+d}}\right)^{1/2}.
\end{equation}

The square root means that the group that acts on the
supersymmetries is a double cover of the duality group $\Gamma$.
The extension is by a symmetry $(-1)^F$ of order two that changes
the sign of all fermions. This symmetry is the {\it square} of the
generator ${\EUI} $ of the center of $SU(4)_{\cal R}$. We avoided
getting an extension of $\Gamma$ by the full center by taking
$\vec{\phi}$ to be invariant under $\Gamma$ rather than invariant
up to sign.

In case $G$ is not simply-laced, this argument needs to be phrased
more carefully.  The duality group is generated by $T:\tau\to
\tau+1$ and $S:\tau\to -1/n_{\frak g}\tau$.  $T$ leaves fixed the
supersymmetries, so $\hat\phi=0$ for $T$, in keeping with \eqref{eq:tokio}.
As for $S$, its action is now % 式 (2.26)
\setcounter{equation}{25}
\begin{equation}\label{eq:xiljo}
\begin{aligned}\vec{\phi} &
\mapsto {\frak R}\cdot\vec{\phi}\\\left(\begin{matrix}\vec{{\eurm m}} &
\vec{{\eurm n}}\end{matrix}\right)&\mapsto \left(\begin{matrix}{\frak R}\cdot
\vec{{\eurm m}} & {\frak R}\cdot\vec{{\eurm
n}}\end{matrix}\right)\left(\begin{matrix}0 & -1/\sqrt n_{\frak g} \\ \sqrt
n_{\frak g} & 0 \\\end{matrix}\right).\end{aligned}
\end{equation}
 Here $\frak R$ is an orthogonal
transformation of $\frak t$ which for simply-laced groups can be
taken to be trivial, and for non-simply-laced groups is described
in \cite{ref47}. The square of $\frak R$ belongs to the Weyl group, so the
square of $S$ acts trivially on the moduli space. $\frak R$ does
not affect the computation of the transformation of the central
charge, and a small computation shows that $n_{\frak g}$ cancels
out, so \eqref{eq:tokio}\ remains valid in the general case.
